\documentclass[11pt]{article}
\usepackage[margin=0.75in]{geometry}
\usepackage[T1]{fontenc}
\usepackage[utf8]{inputenc}
\usepackage{lmodern}
\usepackage{microtype}
\usepackage{array}
\usepackage{booktabs}
\usepackage{longtable}
\usepackage{tabularx}
\usepackage{enumitem}
\usepackage{needspace}
\usepackage{multicol}
\usepackage{xcolor}
\usepackage{hyperref}
\usepackage[most]{tcolorbox}
\usepackage{tikz}
\usetikzlibrary{arrows.meta,calc,positioning,shapes.geometric}

% Semantic color names are retained for compatibility with the weekly files,
% but every value is pure black or white for monochrome printing.
\definecolor{rubricred}{gray}{0}
\definecolor{linkblue}{gray}{0}
\definecolor{softblue}{gray}{1}
\definecolor{softgold}{gray}{1}
\definecolor{softgray}{gray}{1}

\hypersetup{
  hidelinks,
  pdfborder={0 0 0}
}

% PDF timestamps are document source, never build-clock state.  Every
% publishable leaf supplies a strict UTC timestamp through its tracked
% generation-metadata.tex.  Suppressing the automatic creation date and
% trailer ID keeps unchanged source byte-reproducible across rebuilds.
\pdfinfoomitdate=1
\pdftrailerid{}

\setlength{\parindent}{0pt}
\setlength{\parskip}{0.55em}
\setlength{\tabcolsep}{3pt}
\setlist[itemize]{leftmargin=1.25em, itemsep=0.2em}
\renewcommand{\arraystretch}{1.22}

\newcommand{\weektitle}[3]{
  \begin{center}
    {\Large\bfseries #1}\\[0.2em]
    {\large #2}\\[0.2em]
    {\small #3}
  \end{center}
  \vspace{0.5em}
}

\newcommand{\properrow}[4]{\textbf{#1} & #2 & #3 & #4\\}
\newcommand{\properrefs}[1]{\textnormal{(\textit{#1})}}
\newcommand{\sourceurl}[2]{\href{#1}{#2}}
\newcommand{\latin}[1]{\textit{#1}}
\newcommand{\rubric}[1]{\textbf{#1}}

% Generation provenance and the document revision timestamp are stored in each
% document leaf and rendered through these shared commands.  The first model
% contribution renders the timestamp in a canonical provenance block.  A
% mechanically derived companion has no local contribution, so its timestamp
% is rendered once by the universal rights page instead.
\newif\ifAIDocumentRevisionRendered
\def\AIDocumentRevisionTimestampValue{}
\makeatletter
\def\AI@revisiontopdf#1-#2-#3T#4:#5:#6Z\@nil{D:#1#2#3#4#5#6Z}
\newcommand{\AIDocumentRevisionTimestamp}[1]{%
  \gdef\AIDocumentRevisionTimestampValue{#1}%
  \edef\AI@pdfrevisiontimestamp{%
    \expandafter\AI@revisiontopdf#1\@nil}%
  % Metadata files are imported after the first page in some documents, when
  % hyperref has frozen its setup keys.  The pdfTeX primitive remains valid at
  % that point and records the tracked value without consulting the build clock.
  \pdfinfo{/ModDate (\AI@pdfrevisiontimestamp)}%
}
\makeatother
\newcommand{\AIDisplayDocumentRevisionTimestamp}{%
  \ifAIDocumentRevisionRendered\else
    \textbf{Last revised (UTC):} \AIDocumentRevisionTimestampValue\\
    \global\AIDocumentRevisionRenderedtrue
  \fi}
\newcommand{\AIModelContribution}[3]{%
  \AIDisplayDocumentRevisionTimestamp
  \textbf{Model:} \texttt{#1}; \texttt{#2}\\
  \textbf{Agent/runtime:} #3\par}
\newcommand{\AIInheritedGenerationMetadata}[1]{}

% Every detached publication carries the repository's content-license boundary.
% The license covers only project-created contributions, not incorporated source
% material whose rights or public-domain status remain independent.
\newcommand{\TriptychRightsNotice}{%
  \clearpage
  \thispagestyle{empty}
  \null\vfill
  \begingroup
  \centering
  \begin{minipage}{0.82\linewidth}
  \small
  {\large\bfseries Reuse and Rights\par}
  \medskip
  \AIDisplayDocumentRevisionTimestamp
  \medskip
  To the extent the project holds the relevant rights, project-created
  editorial content, annotations, translations, and design are licensed under
  the
  \href{https://creativecommons.org/licenses/by/4.0/}{Creative Commons
  Attribution 4.0 International License}.

  This license does not cover Scripture; liturgical or official texts;
  received prayers or hymns; quotations; fonts; or other third-party material,
  which retain their own status. It creates no exclusive rights in
  public-domain material.

  Changes must be identified. Attribution may identify Triptych as the source,
  but it does not imply Triptych or ecclesiastical approval. The complete
  license map and detailed exclusions are recorded in the \texttt{LICENSE} and
  \texttt{THIRD\_PARTY.md} files distributed with the source.
  \end{minipage}
  \endgroup
  \vfill\null
}
\AtEndDocument{\TriptychRightsNotice}

\newenvironment{properstable}
{\footnotesize\begin{longtable}{>{\raggedright\arraybackslash}p{0.15\linewidth}>{\raggedright\arraybackslash}p{0.25\linewidth}>{\raggedright\arraybackslash}p{0.18\linewidth}>{\raggedright\arraybackslash}p{0.34\linewidth}}
\textbf{Proper} & \textbf{TLM text / reference} & \textbf{Scriptural axis} & \textbf{Connection}\\
\midrule
\endhead}
{\bottomrule\end{longtable}}

\newenvironment{sensestable}
{\vspace{-0.35em}\begingroup\footnotesize\renewcommand{\arraystretch}{1.10}%
\noindent\hspace*{\dimexpr0.04\linewidth-4\tabcolsep\relax}%
\begin{tabular}{>{\bfseries\raggedright\arraybackslash}p{0.15\linewidth}>{\raggedright\arraybackslash}p{\dimexpr0.77\linewidth+4\tabcolsep\relax}}
\textbf{Sense} & \textbf{Synthesis}\\
\midrule}
{\bottomrule\end{tabular}\par\endgroup}

\newenvironment{connectionstable}
{\small\begin{longtable}{>{\bfseries\raggedright\arraybackslash}p{0.18\linewidth}>{\raggedright\arraybackslash}p{0.74\linewidth}}
\toprule
\textbf{Thread} & \textbf{Synthesis}\\
\midrule
\endhead}
{\bottomrule\end{longtable}}

\newenvironment{witnessstable}
{\small\begin{longtable}{>{\raggedright\arraybackslash}p{0.2\linewidth}>{\raggedright\arraybackslash}p{0.34\linewidth}>{\raggedright\arraybackslash}p{0.38\linewidth}}
\toprule
\textbf{Witness} & \textbf{Text / source status} & \textbf{Interpretive contribution}\\
\midrule
\endhead}
{\bottomrule\end{longtable}}

\newcommand{\witnessrow}[3]{#1 & #2 & #3\\}
\newcommand{\sectionguard}{\Needspace{7\baselineskip}}

% Keep the opening survey page compact but unchanged; sustained analysis uses
% a slightly smaller reading size so longer patristic dossiers remain usable.
\newenvironment{deepstudy}
  {\begingroup\small}
  {\par\endgroup}

\newenvironment{threestable}[3]
{\small\begin{longtable}{>{\bfseries\raggedright\arraybackslash}p{0.18\linewidth}>{\raggedright\arraybackslash}p{0.34\linewidth}>{\raggedright\arraybackslash}p{0.40\linewidth}}
\toprule
\textbf{#1} & \textbf{#2} & \textbf{#3}\\
\midrule
\endhead}
{\bottomrule\end{longtable}}

\newenvironment{fourstable}[4]
{\footnotesize\begin{longtable}{>{\bfseries\raggedright\arraybackslash}p{0.15\linewidth}>{\raggedright\arraybackslash}p{0.23\linewidth}>{\raggedright\arraybackslash}p{0.25\linewidth}>{\raggedright\arraybackslash}p{0.27\linewidth}}
\toprule
\textbf{#1} & \textbf{#2} & \textbf{#3} & \textbf{#4}\\
\midrule
\endhead}
{\bottomrule\end{longtable}}

\newtcolorbox{studybox}[1]{
  enhanced,
  breakable,
  colback=white,
  colbacktitle=white,
  coltitle=black,
  frame hidden,
  boxrule=0pt,
  sharp corners,
  borderline west={1.2pt}{0pt}{black},
  left=8pt,
  right=2pt,
  top=2pt,
  bottom=2pt,
  toptitle=0pt,
  bottomtitle=2pt,
  before skip=8pt,
  after skip=8pt,
  title={#1},
  fonttitle=\small\bfseries
}

% Use the titleless form when a surrounding heading and explicit field labels
% already state the block's function.  The left rule supplies hierarchy without
% adding a second prose label that merely narrates the typography.
\newtcolorbox{studyframe}{
  enhanced,
  breakable,
  colback=white,
  frame hidden,
  boxrule=0pt,
  sharp corners,
  borderline west={1.2pt}{0pt}{black},
  left=8pt,
  right=2pt,
  top=4pt,
  bottom=4pt,
  before skip=8pt,
  after skip=8pt
}

\newtcolorbox{massbox}[1]{
  enhanced,
  breakable,
  colback=white,
  colframe=black,
  colbacktitle=white,
  coltitle=black,
  boxrule=0.55pt,
  sharp corners,
  left=7pt,
  right=7pt,
  top=4pt,
  bottom=4pt,
  toptitle=3pt,
  bottomtitle=3pt,
  titlerule=0.45pt,
  before skip=8pt,
  after skip=8pt,
  title={#1},
  fonttitle=\small\bfseries
}

\newtcolorbox{sourcecard}[1]{
  enhanced,
  colback=white,
  colframe=black,
  colbacktitle=white,
  coltitle=black,
  boxrule=0.4pt,
  sharp corners,
  left=5pt,
  right=5pt,
  top=4pt,
  bottom=4pt,
  toptitle=3pt,
  bottomtitle=3pt,
  titlerule=0.4pt,
  before skip=4pt,
  after skip=4pt,
  title={#1},
  fonttitle=\small\bfseries
}

\tikzset{
  guide node/.style={draw=black, fill=white, align=center, inner sep=5pt, font=\small},
  guide emphasis/.style={draw=black, very thick, fill=white, align=center, inner sep=5pt, font=\small},
  guide arrow/.style={-{Stealth[length=2mm]}, thick, draw=black}
}
